
\begin{exercicio}
  Resolva as equações.
  \begin{mathlist}[2]
      \item \frac{x - 2}{x + 3} = 0
      \item \frac{2x + 5}{x - 1} = 3
      \item \frac{5x - 2}{1 - 3x} = -1
      \item \frac{3 - \sfrac{x}{2}}{3x + 8} = \frac{1}{4}
      \item \frac{3x + 5}{4x - 5} = -3
      \item \frac{4 - \sfrac{x}{2}}{4x + 1} = 0
      \item \frac{x}{x^2 - 3x + 2} = 0
      \item \frac{2x^2}{x + 5} = 5
      \item \frac{x^2}{3x - 2} = 2x - 1
      \item \frac{2 \left(x - \sfrac{5}{6}\right) + 1}{5x - 3} = \frac{2}{3}
      \item \frac{x^2 - 26}{x^2 - 9} = 10
      \item \frac{x^2 + 1}{x^2 - 4} = 6
      \item \frac{x^2 + 3x - 10}{x^2 - 2} = 5
      \item \frac{x^2 - 1}{x^2 - 2x + 1} = 4
      \item \frac{x^2 - 4x + 9}{x^2 - 5x + 6} = 3
      \item \frac{2}{x + 1} - \frac{4}{x - 1} = 0
      \item \frac{4}{x + 1} + \frac{1}{x - 1} = \frac{5}{x^2 - 1}
      \item \frac{3}{x + 1} + \frac{2}{x - 1} = 3
      \item \frac{2}{x - 4} + \frac{5}{x - 2} = 3
      \item \frac{6}{x - 3} + \frac{5}{x - 4} = 2
      \item \frac{x}{x - 2} - \frac{3}{x + 2} = 1
      \item \frac{1}{2x + 1} + \frac{1}{3x - 1} = \frac{2}{5}
      \item \frac{3}{x - 2} - \frac{2}{x + 3} = \frac{1}{x}
      \item \frac{2}{x + 1} - \frac{2}{2x - 3} = \frac{3}{x}
      \item \frac{4}{x - 5} - \frac{2}{x - 2} = \frac{12}{(x - 5) (x - 2)}
      \item \frac{4}{5 - 2x} - \frac{3}{x} = 0
  \end{mathlist}    
  \begin{answers}
      \begin{mathlist}[2]
        \item x_1 = 2 
        \item x_1 = 8 
        \item x_1 = \sfrac{1}{2} 
        \item x_1 = \sfrac{4}{5} 
        \item x_1 = \sfrac{2}{3} 
        \item x_1 = 8 
        \item x_1 = 0 
        \item x_1 = - \sfrac{5}{2} \qquad x_2 = 5 
        \item x_1 = \sfrac{2}{5} \qquad x_2 = 1 
        \item x_1 = 1 
        \item x_1 = - \sfrac{8}{3} \qquad x_2 = \sfrac{8}{3} 
        \item x_1 = \sqrt{5} \qquad x_2 = - \sqrt{5} 
        \item x_1 = 0 \qquad x_2 = \sfrac{3}{4} 
        \item x_1 = \sfrac{5}{3} 
        \item x_1 = 1 \qquad x_2 = \sfrac{9}{2} 
        \item x_1 = -3 
        \item x_1 = \sfrac{8}{5} 
        \item x_1 = - \sfrac{1}{3} \qquad x_2 = 2 
        \item x_1 = 3 \qquad x_2 = \sfrac{16}{3} 
        \item x_1 = \sfrac{7}{2} \qquad x_2 = 9 
        \item x_1 = 10 
        \item x_1 = - \sfrac{1}{12} \qquad x_2 = 2 
        \item x_1 = - \sfrac{1}{2} 
        \item x_1 = - \sfrac{9}{4} \qquad x_2 = 1 
        \item \) Não há soluções
        \item x_1 = \sfrac{3}{2} 
    \end{mathlist}
  \end{answers}
\end{exercicio}

\begin{exercicio}
  Resolva as equações.
  \begin{mathlist}[2]
      \item \frac{3x + 4}{1 - 5x} + \frac{3}{2} = \frac{5}{6}
      \item \sqrt{3x + 4} = 8
      \item \sqrt{x + 1} = 2x - 1
      \item \sqrt{2x + 1} = x - 1
      \item \sqrt{x - 3} + x = 9
      \item \sqrt{4 - x} + 2 = 3x
      \item 4 \sqrt{3x - 1} = \frac{2}{3} - 2x
      \item \sqrt{5 - x^2} = 3 - 2x
      \item \sqrt{8x + 25} - 2 = 3 - 4x
      \item \sqrt{4x + 4} - x + 2 = 0
      \item 2x + \sqrt{10 - 4x} = 1
      \item \sqrt{4x + 5} - x = 2x + 4
      \item \sqrt{2x + 5} - 1 = x
      \item 3 + \sqrt{45 - 6x} = x
      \item \sqrt{2x^2 + 7} = 2x - 1
      \item \sqrt{25 - 3x^2} = -x
      \item 2 \sqrt{9x^2 - 7} = 6
      \item \sqrt{x^2 + 3} + x = 5
      \item \sqrt{4x^2 - 1} + 1 = 2x
      \item \sqrt{9x^2 - 2} - 3x = 4
      \item (x + 2)^{2/3} = 9
      \item (5x - 6)^{3/2} = 8
      \item (4x^2 - 13)^{3/5} = 27
      \item x^{1/2} - 13x^{1/4} + 12 = 0
      \item 3 \sqrt{x} - \sqrt[4]{x} = 2
      \item x^{1/3} - x^{1/6} = 2
  \end{mathlist}
  \begin{answers}
      \begin{mathlist}[2]
        \item x_1 = 14 
        \item x_1 = 20 
        \item x_1 = \sfrac{5}{4} 
        \item x_1 = 4 
        \item x_1 = 7 
        \item x_1 = \sfrac{11}{9} 
        \item x_1 = \sfrac{1}{3} 
        \item x_1 = \sfrac{2}{5} 
        \item x_1 = 0 
        \item x_1 = 8 
        \item x_1 = - \sfrac{3}{2} 
        \item x_1 = - \sfrac{11}{9} \qquad x_2 = -1 
        \item x_1 = 2 
        \item x_1 = 6 
        \item x_1 = 3 
        \item x_1 = - \sfrac{5}{2} 
        \item x_1 = - \sfrac{4}{3} \qquad x_2 = \sfrac{4}{3} 
        \item x_1 = \sfrac{11}{5} 
        \item x_1 = \sfrac{1}{2} 
        \item x_1 = - \sfrac{3}{4} 
        \item x_1 = 25 
        \item \) Não há soluções
        \item \) Não há soluções
        \item x_1 = 1 \qquad x_2 = 20736 
        \item x_1 = 1 
        \item x_1 = 64 
    \end{mathlist}
  \end{answers}
\end{exercicio}

